\section{Design-Based Partial-Document Supervision}\label{sec:min-audit}

The local laws become useful when they are checked. C-TreePO treats each
realized tree as a finite population of auditable partial-document units.
Leaves test C1, stored summaries test C2, merge nodes test C3, and optional
post-processing rounds test re-summary drift. A sampling design selects units
with logged positive propensities; an accessible judge or gold oracle scores
the sampled checks; inverse-propensity weighting transports those local
observations back to the finite population of tree units.

The LLM setting adds a form of evidence that classical sketches do not have:
the state itself is readable. A HyperLogLog state is a register array, and a
private mergeable sketch is a noisy bitmap. Here $g(x)$ is natural language.
The audit can therefore combine scored node labels with direct inspection:
does the summary contain the load-bearing claims, qualifications, and boundary
facts? It can also sample counterfactual fibers--pairs of spans or summaries
that look equivalent under $g$--and check whether the trusted oracle agrees.
Those checks do not replace the design-based estimator, but they make the
sufficiency failure mode visible before it reaches the root.

The design-based half of the method starts from DSL's core warning:
gold-standard labels matter because model predictions can look accurate while
inducing biased downstream estimates. The same logic applies inside a
compression tree. Convenient-node inspection measures an unidentified population
of local failures. Randomly sampled tree units with logged inclusion
probabilities give Horvitz--Thompson or H\'ajek-style estimates of the target
finite-population signal. C-TreePO extends the sampled population from
documents and final labels to the leaves and merges that determine whether
the labeled object still represents the original document.

\begin{figure}[t]
\centering
\includegraphics[width=\linewidth]{01_base_sampled_grid.pdf}
\caption{\textbf{Six node-sample sets on the same tree.} Each panel
highlights one possible sampled set of tree units. Highlighted nodes
contribute to the local-law distortion estimate, weighted by their logged
propensities. The root is one eligible tree unit among leaves, stored
summaries, and merge nodes.}
\label{fig:min-audit-sample}
\end{figure}

The audit reports two different quantities. \emph{Violation rates} are
diagnostic: which local law fails, and how often? \emph{Distortion means} are
certificate inputs: how far, in oracle units, is the realized tree from the
projection target where C1, C2, and C3 hold? A deployment should report both.
A low violation rate can still be costly if the failing nodes have large
distortion, and a high violation rate can be less damaging if each failure is
small in the task metric.

The natural design is nested. First sample documents from the target corpus.
Then sample nodes within selected documents. If the audit compares candidate
summaries or pairwise preferences, sample those actions within the selected
node. The logged probability for a labeled unit is the product of the relevant
design probabilities, for example
$p_i=p_d\,p_{u\mid d}$ for a node audit or
$p_i=p_d\,p_{u\mid d}\,p_{a\mid u}$ for an action-level label. Adaptive or
bandit-style sampling is allowed when it has an exploration floor and logged
propensities. Those two records keep the audit tied to the intended
population.

For a realized audit population $\mathcal U$ of size $N$, unit $i$ has
inclusion indicator $Z_i$, logged propensity $\pi_i>0$, and distortion
contribution $Y_i$. The basic Horvitz--Thompson mean is
\[
  \hat\mu_{\mathrm{HT}}
  =
  \frac{1}{N}\sum_{i\in\mathcal U}\frac{Z_i}{\pi_i}Y_i.
\]
The node-level object differs from a root-observed corpus evaluation. In a
root-observed evaluation, every sampled document already has a full-document
target, and the root prediction is scored directly against it.

Local laws are checked relative to the oracle available for the audit. In
root-observed benchmarks, the oracle is estimated from held-out or
sample-split full-document labels. In online node audits, gold-standard labels
can be randomly assigned to leaves and merges with logged propensities. Those
labels play two roles at once: they supervise the tree operators and estimate
the finite-population local-law residual. Calibration to the trusted
social-science target remains a separate term in the certificate.

The design choice is how to combine roots and nodes. Roots calibrate the
target: they say what the full-document expert judgment is. Nodes identify
where compression can lose information before that judgment is made. Under the
oracle/state/local-law assumptions, node-level audits supply cheaper
information about the same target and can be aggregated back to the realized
tree population. When the assumptions fail, the audit reports distortion in
the units of the target.

The certificate has four pieces:
\[
  |G|
  \le
  C_{\mathrm{meth}}\hat{\Delta}_R
  + B_{\mathrm{cal}} + B_{\mathrm{est}} + B_{\mathrm{clip}}.
\]
$\hat{\Delta}_R$ is observed local distortion after $R$ rounds.
$C_{\mathrm{meth}}$ converts oracle distortion into the downstream objective
unit. $B_{\mathrm{cal}}$ covers the gap between the accessible judge and the
trusted target oracle. $B_{\mathrm{est}}$ is sampling error from querying only
some nodes. $B_{\mathrm{clip}}$ records stabilization bias from clipped
weights. More local labels shrink the design-based estimation term, while
calibration error between the judge and the social-science target remains
visible.

Granularity is a methodological choice. If a paragraph, section, or merge span
can be audited against the same oracle-preservation property as the full
document, the researcher can shift expert attention from some roots to cheaper
local units and aggregate them without changing the target. If local laws fail,
the audit estimates the resulting distortion and shows which law failed.

Appendix~\ref{app:v4-dsl-audit} gives the design map, logging schema,
clustered uncertainty calculation, and Lean/proof anchors for the
DSL/IPW certificate layer.
